\section{Testing}

\par Testing was approached with a clear separation in mind: the parts of the system that contain decision logic were tested in isolation, and the parts that talk to the database were tested against a real database. This split keeps the everyday test run fast while still giving confidence that the data layer behaves correctly, since most bugs in a system like this hide either in the business rules or in the gap between the code and the database rather than in the wiring between them.

\subsection{Unit testing the business logic}

\par The first group of tests targets the service and utility layers, where the rules that govern the workflow actually live. These tests do not touch the database or any external service; instead, the repositories that a service depends on are replaced with lightweight fakes that return whatever the test needs. This makes each test fast and predictable, and it means a failing test points directly at a piece of logic rather than at the environment around it.

\par The validation functions were a natural starting point, since they encode the rules a request or a document must satisfy before it is accepted. The tests check, for example, that a request without a template is rejected, that a due date in the past is not allowed, that an uploaded file is refused when it is empty, too large, or has a disallowed extension, and that a comment is rejected when it is too short or consists only of whitespace. The status computation, which decides whether a request is pending, uploaded, or overdue based on its documents and its due date, was tested across the combinations that matter, because this is the kind of logic where an off-by-one mistake is easy to make and hard to notice by hand.

\par A second set of tests covers the access-control rules, which are central to a system that exposes different functionality to citizens, department employees, and administrators. These tests confirm that an administrator can reach any request, that a department member can reach requests belonging to their own department, that the assignee of a request can reach it even from another department, and that an unrelated user is refused. The same approach was used for the rate-limiting rule that prevents a user from posting comments too quickly, and for the rule that only an administrator may generate or delete invitation codes. Verifying these rules with tests is more reassuring than checking them by hand, because an accidental change to the logic would now cause a test to fail rather than silently widening who can see what.

\subsection{Integration testing the data layer}

\par The second group of tests exercises the repository layer against a real PostgreSQL database rather than a fake. The reason for this is simple: the repositories are mostly hand-written SQL, and the only way to know that a query is correct, that a join returns the right columns, or that a transaction rolls back when it should, is to run it against a real database that enforces the same constraints as the production one.

\par To make this practical, each test run starts a throwaway PostgreSQL instance inside a container, using the same database image and the same schema file that the application ships with. The container is created once at the start of the run and removed at the end, and the tables are cleared between individual tests so that each one begins from a known empty state. This setup requires nothing more than a running container engine, which the continuous integration environment already provides, and it avoids the need for a developer to install and maintain a separate database just to run the tests.

\par The most valuable tests in this group concern transactions. When a citizen's request is created, the request itself and the list of documents it expects are written together, and either all of them must be saved or none of them. The tests confirm both halves of this: a successful transaction leaves the request and all its expected documents in the database, while a transaction that fails partway through, whether because of an application error or because the database rejects an invalid value, leaves no trace at all. This matters because a half-written request, one that exists but has no documents attached, would be invisible to manual testing yet would quietly break the workflow for that citizen.

\par The remaining tests in this group verify the more ordinary database operations: that a comment can be written and read back with the author's name correctly joined from the users table, that comments are correctly scoped to a single request, that invitation codes which have already been used are excluded from the active list, and that the lifecycle operations on a request, such as claiming, closing, and cancelling it, change exactly the fields they are meant to.

\par The integration tests proved their worth immediately by exposing a real defect that the unit tests could not have caught. The function responsible for fetching a single comment by its identifier built its query with a placeholder for the comment identifier but never actually passed the identifier to the query. Against a fake repository this code path looked fine, because the fake simply returned a prepared value. Against a real database, however, the query was rejected outright, which meant that the endpoint for retrieving an individual comment failed every time it was called. The fix was small once the test exposed it, but the bug could easily have gone unnoticed on a less-frequently used path, which is precisely why the data layer was tested against a real database rather than a substitute.

\subsection{Continuous integration}

\par Both groups of tests are run automatically on every change pushed to the repository and on every proposed change, through the project's continuous integration pipeline. The fast unit tests and the slower, container-based integration tests run as separate steps, so a failure in either is reported clearly and a broken change is caught before it can be merged. Running the tests automatically, rather than relying on remembering to run them by hand, is what makes them useful as the system grows, since it removes the possibility of a change quietly breaking existing behaviour.
